\begin{statement}{E19}{Proved}{Homology of all original affine twists}
For the original multiplicities $(3,4)$ and every admissible affine
datum, put $p=12l_0-4r_1-3r_2$ as in R5 and $d=|p|$.
Then $p\ne0$ and
\[
 \pi_1(X)=\mathbb Z/d\mathbb Z,
 \qquad
 H_j(X,\mathbb Z)=
 \begin{cases}
  \mathbb Z&j=0,6,\\
  \mathbb Z/d\mathbb Z&j=1,2,3,4,\\
  0&j=5.
 \end{cases}
\]
In particular every $|p|>1$ member is a rational homology
six-sphere, has algebraic dimension one, and has the deformation
and tangent-cohomology conclusions of E11.
\end{statement}
\begin{proof}
We check and extend the gluing calculation of
\cite[Propositions labelled \texttt{mvred}, \texttt{mumford},
and \texttt{matrix}]{Viaclovsky2026Log}. Its section on homology
assumes that the gluing determinant is $\pm1$. We show explicitly
which argument works for every nonzero determinant.

First specify the marking and the collar change. In the notation
of the conventions, the change from the retained lattice to the
ordered lattice $(\delta,a,b,d_0)$ of that paper is
\[
 R=\begin{pmatrix}
 -2&1&0&1\\-6&0&-1&0\\0&-1&0&0\\-1&0&0&0
 \end{pmatrix},\qquad \det R=-1.
\]
Its intertwining identities are $A_0R=RA_1^{\rm retained}$ and
$A_1R=RA_2^{\rm retained}$. With $\psi=d_0^\vee$ one has
$\psi R=-\gamma$. The inverse affine-generator comparison
therefore sends an invariant translation $w$ to $-Rw$.
In particular, for $u_0=d_0+2a+2b$ and
$u_1=d_0+3a+3b$,
\[
 -R\varepsilon=u_0,\qquad
 -R\varepsilon'=u_1-\delta,\qquad
 -R\widehat\delta=-\delta.
\]
Thus the coefficients of $u_0,u_1$ before changing the cusp
reference section are exactly $r_1,r_2$.

A continuous reference section on the punctured base can be
chosen to agree with the section of the Mumford filling on its
annular collar. Here is the extension comparison. The punctured
base retracts onto a graph, so its torus torsor has obstruction
in $H^2(U,\Lambda)=0$. A difference of sections on the cusp
annulus has a winding class in
$H^1(A_\infty,\Lambda)=\Lambda/(A_\infty-I)\Lambda$.
Restriction from graph cocycles is onto: the cocycle
$(v_0,v_1)$ on its two edges restricts to
\[
 -A_1^{-1}v_1-A_\infty v_0
       \pmod{(A_\infty-I)\Lambda}.
\]
Taking $v_0=0$ proves surjectivity. Smooth sections of the Lie
bundle form a fine sheaf, so every such winding class is realized
by a torus section. The remaining difference has a Lie lift,
which extends by a cutoff on the collar. This constructs the
required reference section and identifies the underlying real
torus bundle with the linear flat bundle using that section as
zero.

In the retained peripheral presentation write $c$ for the
primitive common coinvariant and $x,y$ for the finite meridians.
The relations are $x^3=c^{r_1}$, $y^4=c^{r_2}$ and
$xy=c^{l_0}$, with $c$ central. Replacing
$y$ by $c^{-l_0}y$ is the exact change to zero cusp winding.
It gives finite coefficients
\[
 n_0=r_1,\qquad n_1=r_2-4l_0,
 \qquad \Delta=4n_0+3n_1=-p.             \tag{E.19.1}
\]
The remaining winding component and the other affine translations
only change the coefficients $\rho_i$ of $\delta$ in
$v_i=\rho_i\delta+n_i u_i$. They will be retained in the
gluing matrix. Admissibility gives
$\gcd(n_0,3)=\gcd(n_1,4)=1$, so $\Delta$ is odd and is not
divisible by three. In particular it is nonzero.

The actual van Kampen maps kill $\delta,a$ and the cusp meridian.
The relation sending $a$ to $-b$ then kills $b$. The remaining
generators $c,x$ commute and have relations
\[
 -n_0c+3x=0,\qquad -n_1c-4x=0.
\]
The relation matrix has determinant $\Delta$ and first Smith
divisor one, because its entries include $3,4$. This proves the
asserted fundamental group and first homology.

Let $V$ be the complement of the three filling interiors and
$M_i$ their boundaries. Each $H_1(M_i)$ is free of rank three;
$H_1(V)$ is free of rank three; and each reduced filling fibre
$F_i$ has free $H_1$ of rank two. At a finite fibre the last
claim also follows directly from
\[
 H_1(F_i)=\mathbb Z\langle a,d_0,g_i\rangle/
  \mathbb Z(s_i\rho_i a-n_i d_0+m_i g_i),
 \quad(s_0,s_1)=(3,2),\quad(m_0,m_1)=(3,4),
\]
whose relation is primitive by the admissibility just checked.
In the Mayer--Vietoris sequence the degree-one map is therefore
between free rank-nine groups and has finite cokernel
$H_1(X)$. It is injective. This is the only place the source
used $H_1(X)=0$ to eliminate that term; finite first homology
suffices.

The map $H_2(M_i)\to H_2(F_i)$ is surjective. The subsequent
pair map takes a cycle to its intersection with $F_i$, and this
intersection is zero because $m_iF_i$ is principal. The pair
sequence and Poincare--Lefschetz duality thus give
\[
 H_2(X)=\operatorname{coker}\left(
  \bigoplus_i H^3(F_i)\xrightarrow{\sum\tau_i}H_2(V)
 \right),
\]
where $\tau_i$ sends a loop on $F_i$ to its normal-circle torus.
All these identifications are integral.

For clarity the basis calculation in this cokernel is also
checked. The graph spectral sequence gives
\[
0\longrightarrow(\textstyle\bigwedge^2\Lambda)_\Gamma
 \longrightarrow H_2(V)
 \longrightarrow\ker[(A_0-I),(A_1-I)]\longrightarrow0.
\]
The first term is generated freely by $s=[a\wedge b]$:
the exterior-square relations kill
$\delta\wedge a,\delta\wedge b,a\wedge d_0,b\wedge d_0$
and give $\delta\wedge d_0=-6a\wedge b$; the invariant
form $\xi=a^\vee\wedge b^\vee+6d_0^\vee\wedge\delta^\vee$
evaluates to one on the remaining generator. The kernel on the
right is free of rank five. Its basis consists of the two cusp
vectors $-(A_1\lambda,\lambda)$ for $\lambda=\delta,a$ and
the three vectors $(\delta,0),(u_0,0),(0,u_1)$.
The minor in rows $1,3,4,5,8$ of these five columns is $-1$,
so they form a primitive rank-five subgroup and hence the full
integral kernel.
The two cusp vectors are the boundaries of the section disc
swept by its two extending circle actions. The normalization
sequence gives $H^3(F_\infty)=\mathbb Z^2$, so they generate
its entire image. Removing them leaves the free basis
$(s,t,t_0,t_1)$ of a rank-four quotient.

Choose $m_i\alpha_i+n_i\beta_i=1$ and
$z_i=\alpha_i d_0+\beta_i g_i$. The domain basis is
$(a,z_0,a,z_1)$. A normal-circle torus over a lattice loop
$\lambda$ has fibre component
$\mathcal N_i\lambda$, with
$\mathcal N_i=\sum_{j=0}^{m_i-1}A_i^j$, and its integral
against $\xi$ is $-\xi(\lambda,v_i)$. For $g_i$ these
quantities are $v_i$ and zero, since $m_i g_i=v_i$ and
$\xi(v_i,v_i)=0$. Substitution of
\[
 \mathcal N_0a=-\delta,\quad\mathcal N_1a=-2\delta,
 \quad\mathcal N_0d_0=3u_0,\quad\mathcal N_1d_0=4u_1
\]
and $\xi(a,u_0)=2$, $\xi(a,u_1)=3$ gives the actual matrix
\[
 B=\begin{pmatrix}
 -2n_0&-6\rho_0\alpha_0&-3n_1&-6\rho_1\alpha_1\\
 -1&\rho_0\beta_0&2&-\rho_1\beta_1\\
 0&1&0&0\\0&0&0&1
 \end{pmatrix}.                              \tag{E.19.2}
\]
The last two rows remove the second and fourth columns by
unimodular operations. The remaining two-by-two block has
determinant $-\Delta$ and contains $-1$. Consequently the Smith
form of $B$ is $\operatorname{diag}(1,1,1,|\Delta|)$.
This proves $H_2(X)=\mathbb Z/d\mathbb Z$.

The Euler number is two: smooth and multiple torus quotients
contribute zero, and the single hexagonal fibre contributes two.
Rational Poincare duality now gives $b_3=0$, since $b_1=b_2=0$.
Integral duality and the universal coefficient sequence give
$H_5=0$, $H_4=\operatorname{Ext}(H_1,\mathbb Z)$ and
$H_3=\operatorname{Ext}(H_2,\mathbb Z)$, proving the remaining
groups. E1 applies to its torus fibration because $b_2=0$;
E11 already computes its tangent cohomology and full miniversal
line without using this topology.
\end{proof}
